% SPDX-License-Identifier: CC-BY-4.0
% Copyright (c) 2026 Martin Schröder <info@swedishembedded.com>

\section{Evaluation Methodology}\label{sec:eval}

\subsection{What is estimated}

The decision-relevant quantity is the difference, on text the candidate and every
generator of its training data never saw, between the candidate asked under the prompt it is
deployed with and the \emph{base asked under the same prompt}. The unprompted base is
a secondary comparison. The distinction is not pedantic: in the first version of the
system every comparison against the base also changed the system prompt, so the effect of the
prompt was credited to the fine-tune. On the 31-task Jefferson exam the prompt alone moves the
base from 8 to 13 tasks judged right and its invented specifics from 21 to 6 (\cref{tab:exams}).
The exam therefore asks four arms the same tasks, each answered greedily (temperature 0,
at most 512 new tokens): the base under the default prompt, the base under the candidate's
persona prompt (\emph{prompted}), the candidate under the default prompt, and the candidate
under its persona prompt (\emph{deployed}). The primary comparison is fixed in the exam
manifest, before any model is put to it: deployed versus prompted, endpoint ``share of
tasks judged right on the greedy answer''.

\subsection{Family-level splits}\label{sec:split}

The corpora contain the same letters in several editions (about half of the letters of one
edition are printed in another). A split by file or by position therefore leaks: holding out
one print and training on another holds nothing out. The unit of every split is a
\emph{family}: texts that print the same text, in part or in whole, found as the transitive
closure of ``share at least eight distinct runs of eight words, wherever in either
text''. A run held by more than six texts is boilerplate (a formula of the period's
letters) and counts for nothing. The thresholds are set by what two unrelated texts can share by
idiom (a 12-word phrase is five runs and must not merge) against what a reused passage shares
(a 40-word passage is 33 runs).

Three defects in earlier versions of this rule are instructive, because each passed every
test that looked at the split itself:
\begin{itemize}
\item \textbf{A windowed comparison.} Texts were grouped by runs sampled from their first
  700 words. On 3,005 letters, comparing whole texts found 93 groups the windowed rule had
  left apart, 21 of them pairs sharing 50 or more distinct runs: letters printed by one edition
  inside a neighbour whose heading the parser had not cut at. Of the 1,625 files a run had
  been pointed at, 51 shared a passage of eight runs or more with a letter held out for the
  independent exam. Whole-text comparison of all 3,005 letters takes three seconds, so the
  window bought nothing.
\item \textbf{A parser that glued works into letters.} A letter ended only at the next
  letter's heading, so the last letter of a volume swallowed the works and appendices printed
  after it (one ``letter'' held 173,468 words, another 74,501). The overlap rule then
  merged them with thirteen letters, the \emph{Summary View}, the \emph{Manual} and the
  \emph{Notes on the State of Virginia} into one group of 16 files holding 419,188 of the
  1,553,235 words of the corpus (27\%): a group that can never be a meaningful held-out
  unit, and that made one of the monitoring families an outsized, non-letter text. A letter
  now ends at the first book, part, appendix or index heading, a stretch over 15,000 words
  is no letter, and the materials step refuses to write when an exam letter is one text with
  anything it would write.
\item \textbf{An exam that was not clear of the materials.} Of 413 families in the
  independent check, five had at least 90\% of their eight-word runs present in the
  training materials (and 44 had at least 5\%), because the glued ``letters'' carried them.
\end{itemize}

To remove the dependence on any rule being right after the fact, the exam's families are
\emph{reserved first}: a stable hash of the family name picks the families of the final
test and of the dev suite before any task is generated or dataset built, the text of every
edition of each is removed from every source that later stages read, and the exam tasks are
written from the reserved text alone. A candidate is never examined on a family it, or a
release it continues, was trained on. Every derivative record (a paraphrase, a description,
a verbatim passage, an abstention) inherits the family of its source and is confined to that
family's split.

\subsection{Tasks, answers and the judge}

Exam tasks are generated by the same admission rules as training tasks (grounded in the
reserved passage, self-contained, naming their subject). The reference for a task is the
source passage itself. Each answer is judged pointwise and blind to the arm: the judge
sees the task, the reference passage and the answer, and decodes greedily, so one answer
receives one verdict whenever it is judged, and no order among arms can reach a verdict.
The judge's prompt tells it to pass answers consistent with the reference, to fail answers that
contradict it, answer something else, dodge or only say they cannot tell, and to abstain
when it cannot decide; the judge's version number is bumped whenever its prompt, decoding or
parsing changes, so a calibration measured on another reading is not reused.

\paragraph{Calibration and controls.}
A judge is calibrated before its verdicts are read. The easy controls are each task's own
reference as the answer (must pass) and the reference of another family's task as the
answer (must fail). These only measure trivial cases, so the exam adds \emph{hard}
controls: each task answered with the reference of the task from another family whose reference
is closest in vocabulary, wrong by construction, so the share passed is the judge's
false-pass rate on a topical near miss. A judge with pass or fail precision below 0.9 is
untrusted and no claim is made. Qwen3-14B reached precision 1.0 on both on 48 easy controls in
every Jefferson exam; on the Adams exam pass precision was 1.0, fail precision 0.96, abstain
rate 0.02, with two controls misjudged. These are measurements on easy controls only. Agreement
with a person has not yet been measured (\cref{sec:limits}); the tooling exports a blind,
stratified sample (by arm, judge verdict and answer length) for a person to label and then
reports agreement, kappa, false-right and false-wrong rates and length bias.

\paragraph{Bias controls.}
Position bias is excluded by pointwise grading. Verbosity and self-preference
\citep{zheng2023judging,dubois2024lc} are handled less completely. The judge is told that
saying more than the reference is not a fault; length is otherwise not controlled except by
a \emph{like-length test}, which repeats the sign test on the tasks both arms answered within
a quarter of each other's length. The judge and the policy are both Qwen3, so
self-preference cannot be excluded. Deterministic checks override the judge wherever they apply: a recall
answer is decided by whether it states the reference, an executable task by running it, and
a separate code check marks an answer as having \emph{invented specifics} if it states a
number or more than one proper name that the task's source does not hold. This count is
reported beside the judged-right count because the judge is told to credit agreement with
the reference, not to detect extra claims. It partly measures verbosity, since a longer
answer states more specifics; answer lengths are recorded per arm so that this can be
examined.

\subsection{Statistics}

Tasks about one family are not independent, so the family is the unit of resampling. The
primary analysis pairs the arms by task: an exact one-sided sign test over the tasks only
one arm got right, reported with the difference in the share right and its 95\% percentile
bootstrap interval over families (10,000 resamples, seeded, each family contributing its own
tasks). The same test over families, each counted once by which arm got more of its tasks
right, is reported beside it; this is the test the release gate has always used. Secondary
comparisons (candidate under the default prompt, versus the unprompted base, and so on)
are corrected with Holm's step-down procedure~\citep{holm1979}; the primary comparison is
decided on its own. The report estimates the exam's own discordance and intraclass
correlation of the per-task difference (one-way analysis of variance estimator, floored at
zero) and the smallest true difference it could have shown, so that a null can be read as a
bound. Optional extra sampled answers per task give a continuous pass rate that is reported
beside, never in place of, the binary verdict.

\subsection{Power}\label{sec:power}

A family-level sign test needs a clean sweep. \Cref{fig:power} gives the exact power of
the two-sided test at $\alpha=0.05$ as a function of the number $n$ of discordant
families (those in which one arm got strictly more tasks right) and the probability $q$
that the better arm wins a discordant family. With six discordant families the test has
power \PowerSixNine\ even when the better arm wins 90\% of them, and \PowerSixEight\ at
80\%; reaching 80\% power at $q=0.8$ needs about 20 discordant families
(\PowerTwentyEight), and 40 give \PowerFortyEight. Our exams had between two and six
discordant families against the prompted base (\cref{tab:exams}). They could not have
succeeded unless every discordant family was a win and there were at least five of them
(one-sided) or six (two-sided).

\begin{figure}[t]
\centering
\begin{tikzpicture}
\begin{axis}[width=0.8\textwidth, height=6cm, xlabel={discordant families $n$},
  ylabel={power of the two-sided sign test}, xmin=1, xmax=60, ymin=0, ymax=1.02,
  legend pos=south east, legend style={font=\footnotesize}, grid=major, grid style={gray!25},
  tick label style={font=\footnotesize}, label style={font=\small}]
\addplot[thick, blue!60!black] table[x=n, y=q090, col sep=comma] {generated/sign_power.csv};
\addlegendentry{$q=0.9$}
\addplot[thick, green!50!black, dashed] table[x=n, y=q080, col sep=comma] {generated/sign_power.csv};
\addlegendentry{$q=0.8$}
\addplot[thick, orange!80!black, dashdotted] table[x=n, y=q070, col sep=comma] {generated/sign_power.csv};
\addlegendentry{$q=0.7$}
\addplot[thick, red!70!black, dotted] table[x=n, y=q060, col sep=comma] {generated/sign_power.csv};
\addlegendentry{$q=0.6$}
\draw[gray, thick] (axis cs:7,0) -- (axis cs:7,1.02) node[pos=0.5, right, font=\scriptsize, fill=white] {7 families};
\end{axis}
\end{tikzpicture}
\caption{Exact power of the family-level sign test ($\alpha=0.05$, two-sided) against the
number of discordant families $n$ and the per-family win probability $q$ of the better
arm. The vertical line marks the seven families of the pilot-scale exams.}
\label{fig:power}
\end{figure}

A task-level paired test is far more powerful, but its p-value depends on how the
discordant tasks split, which the stored marginals do not reveal. What the marginals do allow
is a bound: with $m$ net extra tasks right, the smallest one-sided exact sign p over all
possible splits is $0.5^m$ (\cref{tab:bounds}). Only the Adams exam (+8 net) could ever have
reached significance at the task level, ignoring clustering.

\paragraph{Sizing the final test by simulation.}
The planned test is the one the exam reports as primary: the deployed arm is better when
the lower end of the family-clustered bootstrap interval of the difference is above zero. We
simulate it: a task's difference is $+1$, $-1$ or $0$; a share \emph{discordance} of tasks
are $\pm1$; the first arm leads by \emph{effect}; and a family's mean difference is the
overall one plus a draw whose variance is the intraclass correlation (ICC) times a task's
variance. With no pilot estimate yet, we assume discordance 0.30 and ICC 0.10, so the
design effect is $1+(8-1)\times0.1=1.7$ for eight tasks per family.
\Cref{tab:power} gives the results. Two points matter. First, the planned 50 families of
eight tasks reach power 0.72, not the 0.8 an earlier estimate from the same assumptions suggested, for
a ten-point effect; it needs a fifteen-point effect to reach 0.97 and has power 0.23 at five
points. Second, the same simulation with 7 families of 5 tasks has a false-positive rate of
0.0725 at zero true effect: a percentile bootstrap over seven clusters is anti-conservative,
so even a nominally significant 7-family interval would deserve scepticism. These simulations
use 400 replicates and 500 bootstrap resamples, so each entry has a Monte Carlo error of about
$\pm0.02$.

\begin{table}[t]
\centering\small
\resizebox{\textwidth}{!}{\input{generated/power_table}}
\caption{Simulated power of the planned paired test (family-clustered bootstrap, lower bound
above zero), assuming a share of 0.30 of tasks discordant and an intraclass correlation of
0.10. ``Realised'' rows use the 39 families that yielded tasks and their mean of six tasks,
since 50 families were reserved but only 39 produced admitted tasks (244 tasks).}
\label{tab:power}
\end{table}

\paragraph{Procedure: pilot, simulate, freeze, open once.}
A pilot on the first 25 families by name (an arbitrary fixed order) estimates discordance and
ICC; the planned test is simulated with those; the size is frozen; a dev suite
(50 families of four tasks reserved beside the final test) chooses the checkpoint by a fixed
rule (at most 8\% invented specifics, then the best answers within a standard error, then
the loss on the dev families' own text); and the final test is run in full once, after one
checkpoint is chosen (the pilot's 25 families are a subset of it). The frozen exam is pinned by a ledger: the same name with other
content is refused, so every model is examined on the same tasks.

\subsection{Retention and the anchor suite}

The anchor suite guards against a persona fine-tune eroding general behaviour. It has 110
items graded by code. The gate computes the drop in accuracy of the candidate against the
base and fails when it exceeds 0.02. This test has two problems that no amount of
grading care removes. First, it is a point estimate on a paired binary outcome: with 17
discordant items the standard deviation of the paired drop under no change is
$\sqrt{17}/110=\AnchorSD$, so a bound of 0.02 sits about half a standard deviation from zero
and is exceeded in about 31\% of draws by chance. Second, to see a true two-point drop with
80\% power at the discordance observed (17 of 110, 0.155) the normal approximation
$n=(z_{\alpha/2}+z_\beta)^2 d/\delta^2$ gives about 3,000 paired items. The gate now
reports the paired bootstrap interval of the drop and the number of items needed beside the point
estimate; the 0.02 bound is unchanged because changing a bound is a change of rule, not a
report.

\begin{table}[t]
\centering\small
\resizebox{\textwidth}{!}{\input{generated/anchor_table}}
\caption{The anchor suite on three candidates. ``Drop'' is the base's accuracy minus the
candidate's; $p$ is the exact two-sided sign test on the discordant items. The bound is
0.02; the first two rows failed it, with $p=0.63$ and $p=0.45$.}
\label{tab:anchor}
\end{table}
